\section{Project Domain Analysis}

\subsection{Domain Description}

Electrical power distribution systems are responsible for delivering
electricity from utility providers to consumers while monitoring energy
consumption for billing purposes. Traditional systems rely on centralized
metering infrastructure and manual verification of billing records. The
proposed project introduces a decentralized approach that combines IoT
smart meters, blockchain technology, and machine learning to provide
transparent prepaid energy management, secure transaction records, and
real-time monitoring.

\subsection{Existing System}

In conventional power distribution systems, smart meters or electricity
meters send consumption data to a centralized utility database. Billing
is generally performed using postpaid models, where consumers receive
their bills after energy has been consumed. Consumers have limited
visibility into how charges are calculated, and the centralized
architecture creates a single point of failure for data storage and
billing operations.

\subsection{Proposed System}

The proposed Decentralized Smart Power Grid continuously receives
electricity consumption data from ESP32-based smart meter nodes. The
backend validates the incoming telemetry and forwards billing requests
to blockchain smart contracts that automatically deduct prepaid credit.
Machine learning models analyze consumption patterns to estimate
remaining credit duration and detect possible power theft or abnormal
usage. The processed information is stored securely and displayed on a
web dashboard, allowing both consumers and administrators to monitor
energy usage, account balance, billing history, and system alerts in
real time.